\documentclass[11pt]{amsart}
\usepackage{amsmath,amssymb,amsthm}
\usepackage[margin=1.05in]{geometry}
\raggedbottom
\usepackage[hidelinks]{hyperref}
\usepackage{url}
\newtheorem{theorem}{Theorem}[section]
\newtheorem{proposition}[theorem]{Proposition}
\newtheorem{lemma}[theorem]{Lemma}
\theoremstyle{remark}
\newtheorem{remark}[theorem]{Remark}
\newcommand{\Z}{\mathbb{Z}}
\newcommand{\Q}{\mathbb{Q}}
\newcommand{\OH}{\mathcal{O}_H}
\newcommand{\pp}{\mathfrak{p}}
\newcommand{\PP}{\mathfrak{P}}
\newcommand{\Ns}{N}
\newcommand{\Tr}{\operatorname{Tr}}
\newcommand{\coker}{\operatorname{coker}}
\newcommand{\Aloc}{\mathcal{A}^{\mathrm{aff},\mathrm{loc}}}

\title[Integral equivariant K-theory of the inert model]{Integral equivariant K-theory of the inert local model:\\ the ray-class entanglement and exact torsion}
\author{Ruqing Chen}
\address{GUT Geoservice Inc., Montreal, Canada}
\email{ruqing@hotmail.com}
\thanks{Paper XXXVIII of the series on localized Bost--Connes systems. The framework is
that of \cite{Main}, cited by DOI; Papers XXX, XXXI and XXXII of the series are referred
to by numeral, and every fact used is restated in place. Search and Smith-form scripts in
\texttt{code/xxxviii-inert/} of the repository.}
\date{6 September 2026}

\begin{document}

\begin{abstract}
We carry out the first integral case of Paper XXXII's Problem (i): the equivariant
K-theory of the inert local model at $K=\Q(\sqrt{-5})$, $H=\Q(i,\sqrt5)$,
$\pp_3=(3,1+\sqrt{-5})$ inert in $H/K$, $C=\mathrm{Gal}(H/K)\cong\Z/2$. The ring
$\OH=\Z[i]\otimes\Z[\varphi]$ is a \emph{free} $\Z[C]$-module of rank two (a Tate
cohomology computation), so the entire exterior tower carries explicit integral
$C$-structure: $\Lambda^0=\Z_+$, $\Lambda^1\cong\Z[C]^2$,
$\Lambda^2\cong\Z_-^2\oplus\Z[C]^2$, $\Lambda^3\cong\Z[C]^2$, $\Lambda^4=\Z_+$. A finite
search produces the explicit generator $\Pi=-3-i+2\varphi+i\varphi$ of a prime above
$3$, with $\Ns_{H/\Q}\Pi=+9$ and $\Tr_{H/\Q}\Pi=-8$, and Smith normal forms of
$9\cdot1-\Lambda^r\Pi$ evaluate the intermediate layers exactly (prime-to-$3$
regularization): $\Z/5\oplus\Z/5\oplus\Z/61$ at $r=1$;
$\Z/2\oplus\Z/2\oplus\Z/5\oplus\Z/7\oplus\Z/11$ at $r=2$, containing the sign-twisted
$\Z_-^2$-block; $\Z/29$ at $r=3$. Assembling with the top and bottom layers gives
\[
K_0\supseteq\Z^2\oplus(\Z/8)^2\oplus\bigl[\Z/2^{\oplus2}\oplus\Z/5\oplus\Z/7\oplus\Z/11\bigr],
\qquad
K_1\supseteq\Z^2\oplus\bigl[\Z/5^{\oplus2}\oplus\Z/61\bigr]\oplus\Z/29,
\]
where $(\Z/8)^2=\Z/(q^h{-}1)\otimes R(\Z/2)$: the unit class is split by the spectral
projections $\tfrac{1\pm U_\sigma}2$ into a pair of classes of exact order $8$. The
entanglement law this computation exhibits: the class group doubles or sign-twists
exactly the \emph{non-induced} strata, while free $\Z[C]$-blocks count once by Shapiro
--- the class group does not change the orders of the middle layers, only their
representation labels. Two honest gaps close the paper: the $u_\sigma$-twisted
bookkeeping of those labels, and the extension data in degree four. A structural
by-product is recorded: the $\sigma$-fixed locus of the inert model is the bottom
$h$-step system of Paper XXXI --- null in measure, nonempty in topology --- so Paper
XXXII's Problem (ii) is a direct summand of Problem (i), not its sequel.
\end{abstract}

\maketitle

\section{The explicit generator and the equivariant exterior tower}\label{sec:tower}

The discriminants $-4$ and $5$ are coprime, so $\OH=\Z[i]\otimes\Z[\varphi]$,
$\varphi=\tfrac{1+\sqrt5}2$, with $\sigma$ (the generator of $C$, fixing
$\sqrt{-5}=i\sqrt5$) acting by $i\mapsto-i$, $\varphi\mapsto1-\varphi$.

\begin{lemma}\label{lem:free}
$\Z[\varphi]\cong\Z[C]$ as a $C$-module ($\hat H^0=\hat H^1=0$: the norm has image $\Z$
and $\ker N=\mathrm{im}(1-\sigma)=\Z\sqrt5$), and $\Z[i]\cong\Z_+\oplus\Z_-$; hence
$\OH\cong\Z[C]^2$ is free, and the integral exterior tower is
\[
\Lambda^0=\Z_+,\quad\Lambda^1\cong\Z[C]^2,\quad
\Lambda^2\cong\Z_-^{\,2}\oplus\Z[C]^2,\quad\Lambda^3\cong\Z[C]^2,\quad\Lambda^4=\Z_+,
\]
the $\Z_-$'s in $\Lambda^2$ coming from $\Lambda^2_\Z(\Z[C])$
($\sigma(1\wedge\sigma)=-(1\wedge\sigma)$), and $\det\sigma=+1$ on $\Lambda^4$.
\end{lemma}

\begin{proposition}[The generator]\label{prop:pi}
$\Pi=-3-i+2\varphi+i\varphi\in\OH$ satisfies $\Ns_{H/\Q}\Pi=+9$ and
$\Tr_{H/\Q}\Pi=-8$; its multiplication matrix $A$ on the basis $(1,i,\varphi,i\varphi)$
has $\det A=+9$, so $(\Pi)=:\PP_1$ is a prime of norm $9$ above $3$, and the inert
model of Paper XXXII at $\PP_1$ is defined with scaling $\Pi$.
\end{proposition}

\section{Exact evaluation of the intermediate layers}\label{sec:layers}

The tower transfers invert only the prime $3$ (they are built from the $\Pi$-adjugates,
of determinant a power of $\Ns\Pi=9$), so the regularized cokernel of $1-\alpha_*$,
$\alpha_*=9^{-1}\Lambda^r\Pi$, is the prime-to-$3$ part of the integral cokernel of
$9\cdot1-\Lambda^r A$, computed by Smith normal form:

\begin{theorem}\label{thm:layers}
With $T_r$ the regularized cokernel at layer $r$:
\[
\begin{array}{lll}
r=1: & \det(9-\Lambda^1A)=13725=3^2\!\cdot\!5^2\!\cdot\!61, &
T_1\cong\Z/5\oplus\Z/5\oplus\Z/61;\\[2pt]
r=2: & \det(9-\Lambda^2A)=-1122660, &
T_2\cong\Z/2\oplus\Z/2\oplus\Z/5\oplus\Z/7\oplus\Z/11;\\[2pt]
r=3: & \det(9-\Lambda^3A)=21141=3^6\!\cdot\!29, &
T_3\cong\Z/29,
\end{array}
\]
$T_2$ containing the contribution of the sign-twisted $\Z_-^{\,2}$-block of
Lemma~\ref{lem:free}. At the ends, $T_0=\Z/(9-1)=\Z/8$ generated by the unit class, and
$M_4=\Z$ with $\alpha_*=\Ns\Pi/9=+1$.
\end{theorem}

\section{The K-theory of the crossed product and the entanglement law}\label{sec:ktheory}

\begin{theorem}\label{thm:main}
For the inert local model at $\PP_1$ with its $\Z/2$-action,
\[
K_0\bigl(\partial\Aloc_{\PP_1}\rtimes\Z/2\bigr)\ \supseteq\
\Z^2\ \oplus\ (\Z/8)^2\ \oplus\
\bigl[\Z/2\oplus\Z/2\oplus\Z/5\oplus\Z/7\oplus\Z/11\bigr],
\]
\[
K_1\bigl(\partial\Aloc_{\PP_1}\rtimes\Z/2\bigr)\ \supseteq\
\Z^2\ \oplus\ \bigl[\Z/5\oplus\Z/5\oplus\Z/61\bigr]\ \oplus\ \Z/29 .
\]
Here $(\Z/8)^2=T_0\otimes R(\Z/2)$: \emph{the unit class is split by the spectral
projections $\tfrac{1\pm U_\sigma}2$ of the implementing unitary into two classes of
exact order $8$}, one per character of the class group; and the two $\Z$'s are the
$R(\Z/2)$-doubling of the trivially-acted top layer.
\end{theorem}

\begin{remark}[The entanglement law]\label{rem:law}
The class group entangles with the $q$-torsion precisely on the \emph{non-induced}
strata: layers with trivial or sign action ($r=0$, $r=4$, and the $\Z_-^{\,2}$-block of
$\Lambda^2$) are doubled or sign-twisted by $R(\Z/2)$, while the free $\Z[C]$-blocks
contribute once, by Shapiro's lemma. The class group does not change the orders of the
middle layers --- those are the Smith invariants of Theorem~\ref{thm:layers} --- it
changes their representation labels.
\end{remark}

\begin{remark}[Problem (ii) is a summand of Problem (i)]\label{rem:fixed}
The $\sigma$-fixed locus of $\mathcal O_{\PP_1}$ is $\mathcal O_{\pp_3}$: null for the
Gibbs measure, nonempty in topology --- and K-theory sees topology. The fixed-point
contribution to the equivariant theory is the bottom $h$-step system of Paper XXXI
re-entering as the fixed subgroupoid: the fixed-point-versus-crossed-product question of
Paper XXXII is a direct summand of the present computation, not its sequel.
\end{remark}

\section{Assessment}\label{sec:assessment}

Two gaps remain and are recorded as such. \emph{(a) Character bookkeeping:} the scaling
is equivariant only up to the descent cocycle $u_\sigma=\sigma\Pi/\Pi$ of Paper XXXII,
so the middle layers are twisted $R(\Z/2)$-modules; their orders are fixed by
Theorem~\ref{thm:layers}, but the assignment of each cyclic summand to the trivial or
sign character awaits the twisted bookkeeping. \emph{(b) Extension data:} in degree four
the Pimsner--Voiculescu assembly determines $K_*$ up to group extensions, as in Paper
XXX. Neither gap touches the displayed subgroups, the split unit class, or the
entanglement law.

\begin{thebibliography}{9}
\bibitem{Main} R. Chen, \emph{KMS state uniqueness and phase transitions in localized Bost--Connes systems}, preprint, 2026, \newline\url{https://doi.org/10.5281/zenodo.22152101}.
\bibitem{CR} C. W. Curtis, I. Reiner, \emph{Methods of Representation Theory I}, Wiley, 1981.
\end{thebibliography}

\end{document}
